
% ==================================================
% 1. INTRODUCTION
% ==================================================

\section{Introduction}

\PAChange{The purpose of this Vision Document is to
identify and describe the high-level needs, scope,
and main features of the Medication Safety Assistant.
It explains the problems faced by target users,
their needs, and the benefits of the proposed system.}

\PAChange{The Medication Safety Assistant is a
responsive Web research prototype focused on
Vietnamese medication information. It supports
medication intake, identification, safety checks,
simulated Healthcare Professional (HCP) review,
medication reminders, adherence tracking, and
patient health summaries. It does not independently
diagnose, prescribe, change treatment, or replace
healthcare professionals.}

\PAChange{This revised document presents the product
positioning, stakeholders and users, user environments,
alternatives and competitors, product features, and
non-functional requirements. Detailed user--system
interactions are described separately in the
Use-case Specifications. Planned features are
identified outside the current prototype scope.}


% ==================================================
% 1.1 REFERENCES
% ==================================================

\subsection{References}

\begin{longtable}{|
    L{0.06\textwidth}|
    L{0.34\textwidth}|
    L{0.20\textwidth}|
    L{0.30\textwidth}|
}
\hline

\textbf{No.} &
\textbf{Document / Source} &
\textbf{Date / Organization} &
\textbf{Location}
\\ \hline

\endfirsthead

\hline
\textbf{No.} &
\textbf{Document / Source} &
\textbf{Date / Organization} &
\textbf{Location}
\\ \hline

\endhead

1
& Medication Safety Assistant Project Proposal
& October 2026, Group 02
& Project repository, PA0 proposal
\\ \hline

2
& Medication Safety Assistant Software
Development Plan, Version 1.0
& October 2026, Group 02
& Project repository, PA1 Project Plan
\\ \hline

3
& Medication Without Harm: Policy Brief
& 7 March 2024,
World Health Organization
& \href{https://www.who.int/publications/i/item/9789240062764}
{WHO publication page}
\\ \hline

4
& Publication of Drug Registration Data in Vietnam
& Accessed 16 October 2026,
Drug Administration of Vietnam
& \href{https://dichvucong.dav.gov.vn/congbothuoc/index}
{DAV registration portal}
\\ \hline

5
& DDInter 2.0: An Enhanced Drug
Interaction Resource
& 2025, Tian et al.,
Nucleic Acids Research
& \href{https://doi.org/10.1093/nar/gkae726}
{DOI: 10.1093/nar/gkae726}
\\ \hline

6
& Web Content Accessibility Guidelines
(WCAG) 2.2
& W3C
& \href{https://www.w3.org/TR/WCAG22/}
{W3C WCAG 2.2}
\\ \hline

\PAChange{7}
& \PAChange{Drugs.com Drug Interaction Checker}
& \PAChange{Accessed 01 November 2026,
Drugs.com}
& \PAChange{\href{https://www.drugs.com/drug_interactions.html}
{Official website}}
\\ \hline

\PAChange{8}
& \PAChange{Medscape Drug Interaction Checker}
& \PAChange{Accessed 01 November 2026,
Medscape}
& \PAChange{\href{https://reference.medscape.com/drug-interactionchecker}
{Official website}}
\\ \hline

\PAChange{9}
& \PAChange{How Can I See Possible
Med Interactions?}
& \PAChange{Accessed 01 November 2026,
Medisafe Help Center}
& \PAChange{\href{https://help-center.medisafe.com/en/articles/8097955-how-can-i-see-possible-med-interactions}
{Medisafe documentation}}
\\ \hline

\end{longtable}
